\begin{table}[ht]
\caption{\textbf{\aops example.} The negative example is retrieved by BM25, and it uses stars and bars to solve the problem, whereas the query and gold document use the Chicken McNugget Theorem.}
\centering
% \resizebox{\textwidth}{!}{
\begin{tabular}{p{13cm}}
\toprule
% \hspace{5.5cm} \sustainable \\
% \midrule
\textit{\textbf{Query} (from 2015 AMC 10B Problem 15)} \\
\midrule
The town of Hamlet has $3$ people for each horse, $4$ sheep for each cow, and $3$ ducks for each person. Which of the following could not possibly be the total number of people, horses, sheep, cows, and ducks in Hamlet? \\
$\textbf{(A) }41\qquad\textbf{(B) }47\qquad\textbf{(C) }59\qquad\textbf{(D) }61\qquad\textbf{(E) }66$ \\
\midrule
\textit{\textbf{Chain-of-thought reasoning to find documents}} \\
\midrule
We can use the \href{https://artofproblemsolving.com/wiki/index.php/Chicken_McNugget_Theorem}{Chicken McNugget Theorem} to solve this problem. We can find other solutions that also apply this theorem.\\
\midrule
\textit{\textbf{Example positive document}} \\
\midrule
Find the sum of all positive integers $n$ such that, given an unlimited supply of stamps of denominations $5,n,$ and $n+1$ cents, $91$ cents is the greatest postage that cannot be formed.

By the Chicken McNugget theorem, the least possible value of $n$ such that $91$ cents cannot be formed satisfies $5n - (5 + n) = 91 \implies n = 24$, so $n$ must be at least $24$.

For a value of $n$ to work, we must not only be unable to form the value $91$, but we must also be able to form the values $92$ through $96$, as with these five values, we can form any value greater than $96$ by using additional $5$ cent stamps.

Notice that we must form the value $96$ without forming the value $91$. If we use any $5$ cent stamps when forming $96$, we could simply remove one to get $91$. This means that we must obtain the value $96$ using only stamps of denominations $n$ and $n+1$.

Recalling that $n \geq 24$, we can easily figure out the working $(n,n+1)$ pairs that can used to obtain $96$, as we can use at most $\frac{96}{24}=4$ stamps without going over. The potential sets are $(24, 25), (31, 32), (32, 33), (47, 48), (48, 49), (95, 96)$, and $(96, 97)$.

The last two obviously do not work, since they are too large to form the values $92$ through $94$, and by a little testing, only $(24, 25)$ and $(47, 48)$ can form the necessary values, so $n \in \{24, 47\}$. $24 + 47 = \boxed{71}$.\\
\midrule
\textit{\textbf{Example negative document}} \\
\midrule
Alice has $24$ apples. In how many ways can she share them with Becky and Chris so that each of the three people has at least two apples?

$\textbf{(A) }105\qquad\textbf{(B) }114\qquad\textbf{(C) }190\qquad\textbf{(D) }210\qquad\textbf{(E) }380$

Note: This solution uses the non-negative version for stars and bars. A solution using the positive version of stars is similar (first removing an apple from each person instead of 2).

This method uses the counting method of stars and bars (non-negative version). Since each person must have at least $2$ apples, we can remove $2*3$ apples from the total that need to be sorted. With the remaining $18$ apples, we can use stars and bars to determine the number of possibilities. Assume there are $18$ stars in a row, and $2$ bars, which will be placed to separate the stars into groups of $3$. In total, there are $18$ spaces for stars $+ 2$ spaces for bars, for a total of $20$ spaces. We can now do $20 \choose 2$. This is because if we choose distinct $2$ spots for the bars to be placed, each combo of $3$ groups will be different, and all apples will add up to $18$. We can also do this because the apples are indistinguishable. $20 \choose 2$ is $190$, therefore the answer is $\boxed{\textbf{(C) }190}$. 
\\
\bottomrule
\end{tabular}
% }
\label{tab:example_math}
\end{table}